\subsection{Reproducing FastGA's order of pairs}\label{sec:impl:order}

\paragraph{Why order matters.} FastGA's alignment loop works through a
contig pair's chains and, after each call to the aligner, continues from
the end position the aligner reached. When a contig is aligned with itself
on the same strand and a chain lies on the main diagonal, the loop skips
the aligner, which would only align the sequence with itself, and resets
the start and end in the first sequence but not in the second. It then
continues from the old end in the second sequence, which belongs to the
previous aligner call, usually for a different contig pair. If that value
lies ahead of the loop's position, the loop jumps to it and can skip real
alignments. Our code copies the loop, so the result for such a pair
depended on which pair a rank aligned just before it, and so on the rank
count (Section~\ref{sec:results:human}).

\paragraph{FastGA's order.} FastGA numbers the contigs by length, longest
first, and stores that order in its index. It cuts the contigs, in this
order, into at most $t$ parts of about equal length, and each part into at
most $t$ consecutive ranges of about equal seed count, one per thread. A
thread aligns its range anchor by anchor, and each anchor's partners in the
same order. Each range starts with no previous aligner call. The value a
self pair $(c,c)$ starts from is therefore the end position left by the
last aligner call before it in its range: usually that of the pair with
$c$'s nearest earlier partner that reached the aligner, otherwise the last
call for an earlier contig in the range.

\paragraph{Supplying the value.} A pair of two different contigs never
reads a leftover value, so its final end position is the same on any rank.
Each rank reads FastGA's contig order from the index, recomputes FastGA's
parts and ranges for the thread count the reference used, and learns from
one reduction which pairs have seeds. It then aligns its same-strand pairs
in a new order: pairs whose partner comes earlier than the anchor, nearest
first, then the other pairs, then the self pairs. After each
earlier-partner pair it sends the pair's final end position, and whether
the pair reached the aligner, to the rank that owns the anchor's self pair:
one small message, received while that rank keeps aligning. A self pair
starts once the answers it needs have arrived; until then it is set aside
and run after the rank's other pairs. The rare value that comes from an
earlier contig is found at the end with two reductions. Opposite-strand
pairs never read the value and are aligned as before.

Where a range's first aligner call is a self pair's own, FastGA reads a
variable it has not set. A build of FastGA that prints the value showed
32,764 to 32,767 for the first range of a run and 0 for the others, on all
four datasets; we use 32,767 and 0.

\paragraph{Cost and limits.} The change needs FastGA's index file beside
the genome database and a map of one byte per pair of contigs, 2\,MB on
hsapiens90.chr6. At 64 ranks, ranks send about 125,000 messages on that
dataset and 1,900 on primates14.chr6. The loop itself is unchanged apart
from setting and reporting the value. Which rank owns a pair can still be
chosen freely, but a self pair cannot start before its nearest earlier
partners have been aligned, wherever they are, which constrains schedules
that reorder pairs within a rank (Section~\ref{sec:disc:future}). An
environment variable restores the original order.
